
\subsection{Test di sistema}
(Da RFF-14 fino a RFF-20) (FYI:il tracciamento c'è sul docs, ma qua non lo metto per adesso, dato che ID non sono ancora giuste)
\begin{tabularx}{\textwidth}{|c|X|c|}
    \hline
    \textbf{ID} & \textbf{Descrizione} & \textbf{Stato}\\
    \hline
    \label{ts:ts} T.S. & Verificare che l’utente possa salvare un test & S\\
    \hline
    \label{ts:ts} T.S. & 
    Verificare che l’utente possa visualizzare la lista dei test salvati, in particolare per ogni test deve visualizzare:
    \begin{itemize}
        \item Il nome del test
        \item Il nome del dataset utilizzato
        \item La data di esecuzione del test
    \end{itemize} 
    E nel caso se non ci fosse presente nessun test salvato l’utente verrà notificato dal sistema
    & S\\
    \hline
    \label{ts:ts} T.S. & 
    Verificare che l’utente possa modificare un test salvato, deve poter quindi:  
    \begin{itemize}
        \item Rinominare un test salvato
        \item eliminare un test salvato
    \end{itemize} 
    & S\\
    \hline
    \label{ts:ts} T.S. & Verificare che l’utente possa caricare un test salvato & S\\
    \hline
    \label{ts:ts} T.S. & Verificare che il sistema possa individuare che il nome di un test salvato inserito dall’utente è invalido & S\\
    \hline
    \label{ts:ts} T.S. & 
    Verificare che l'utente possa visualizzare i risultati del test, deve poter visualizzare quindi: 
    \begin{itemize}
        \item i risultati da uno specificato numero di pagina che vi appartengono
        \item il numero di pagina massima dei risultati
    \end{itemize} 
    & S\\
    \hline
    \label{ts:ts} T.S. & Verificare che il sistema possa individuare che il numero di pagina inserita dall’utente è invalida & S\\
    \hline
    \label{ts:ts} T.S. & Verificare che l’utente possa visualizzare la lista degli LLM salvati, e nel caso se non ci fosse presente nessun LLM salvato l’utente verrà notificato dal sistema & S\\
    \hline
    \label{ts:ts} T.S. & 
    Verificare che l’utente possa modificare un LLM salvato, deve poter quindi: 
    \begin{itemize}
        \item Eliminare un LLM salvato
        \item Rinominare un LLM salvato
    \end{itemize} 
    & S\\
    \hline
    \label{ts:ts} T.S. & Verificare che l’utente possa salvare un nuovo LLM & S\\
    \hline
    \label{ts:ts} T.S. & Verificare che il sistema possa determinare che il nome di un LLM salvato inserito dall’utente è invalido & S\\
    \hline
    \label{ts:ts} T.S. & 
    Verificare che l’utente possa visualizzare in dettaglio un LLM salvato, quindi deve poter visualizzare: 
    \begin{itemize}
        \item il nome dell LLM
        \item la data di ultima modifica agli LLM
        \item l’URL associato
        \item Le chiavi associate
    \end{itemize} 
    & S\\
    \hline
    \label{ts:ts} T.S. & Verificare che il sistema possa individuare che un URL inserito dall’utente è invalido & S\\
    \hline
    \label{ts:ts} T.S. & 
    Verificare che l’utente possa modificare un LLM salvato, quindi deve poter modificare:
    \begin{itemize}
        \item il nome dell LLM
        \item l’URL associato
        \item Le chiavi associate
    \end{itemize} 
    & S\\
    \hline
    \label{ts:ts} T.S. & Verificare che il sistema possa determinare che le chiavi associate inserite dall’utente sono invalide & S\\
    \hline
    \label{ts:ts} T.S. & 
    Verificare che l’utente possa gestire gli insiemi di coppie chiave-valore associate alla creazione delle richieste HTTP verso l’LLM, in particolare deve poter: 
    \begin{itemize}
        \item Visualizzare le coppie chiave-valore
        \item Aggiungere una coppia chiave-valore
        \item Modificare una coppia chiave-valore
    \end{itemize} 
    & S\\
    \hline
    \label{ts:ts} T.S. & Verificare che il sistema possa determinare che una coppia-valore inserita dall’utente è invalido & S\\
    \hline
    \label{ts:ts} T.S. & Verificare che il sistema possa possa determinare che il file JSON selezionato dall’utente è invalido & S\\
    \hline
    \label{ts:ts} T.S. & Verificare che il sistema possa determinare quali test salvati possono essere confrontati con il test caricato & S\\
    \hline
\end{tabularx}


